\chapter{Kubernetes et Helm}
\label{ch:k8s}

\section{Kubernetes en production (k3s)}

\subsection{Principe}

\textbf{Kubernetes} orchestre des conteneurs sur un cluster : scheduling, réseau, stockage, auto-healing via sondes. En production, le projet utilise \textbf{k3s} sur un \textbf{VPS dédié} (séparé de tout serveur Coolify client) : distribution légère, adaptée à un nœud unique, avec ServiceLB pour exposer ingress-nginx sur les ports~80 et~443.

\subsection{Installation}

k3s est installé avec Traefik désactivé (ingress-nginx est géré par Terraform) :

\begin{lstlisting}[style=bash]
curl -sfL https://get.k3s.io | sh -s - - --disable traefik --write-kubeconfig-mode 644
\end{lstlisting}

Le guide opérationnel \texttt{docs/vps-setup.md} détaille DNS, pare-feu, secrets GitHub et bootstrap.

\section{Namespaces}

Terraform crée les namespaces (\texttt{kubernetes.tf}) :
\begin{itemize}
  \item \texttt{ingress-nginx} — contrôleur Ingress ;
  \item \texttt{cert-manager} — certificats TLS Let's Encrypt ;
  \item \texttt{monitoring} — kube-prometheus-stack ;
  \item \texttt{pfa-stock} — application métier ;
  \item \texttt{pgadmin} — administration PostgreSQL ;
  \item \texttt{kubernetes-dashboard} — visualisation du cluster.
\end{itemize}

\section{Chart Helm \texttt{pfa-stock}}

\subsection{Métadonnées}

\texttt{helm/pfa-stock/Chart.yaml} : chart version 1.0.0, application version 1.0.0.

\subsection{Values}

\begin{itemize}
  \item \texttt{values.yaml} — valeurs par défaut (développement / lint) ;
  \item \texttt{values-prod.yaml} — production : sous-domaines \texttt{pfa.elyes.dev} et \texttt{api.pfa.elyes.dev}, TLS via cert-manager, \texttt{imagePullSecrets} pour GHCR.
\end{itemize}

Les tags d'images et les secrets sensibles sont injectés par le workflow \texttt{deploy.yml} (\texttt{--set}, \texttt{--set-string}).

\subsection{Ressources déployées}

\begin{table}[htbp]
  \centering
  \caption{Ressources du chart \texttt{pfa-stock}}
  \label{tab:helm-resources}
  \small
  \begin{tabular}{@{}llp{7.5cm}@{}}
    \toprule
    \textbf{Template} & \textbf{Kind} & \textbf{Détail} \\
    \midrule
    \texttt{postgres.yaml} & StatefulSet + Service & Volume \texttt{1Gi}, image postgres:16-alpine \\
    \texttt{api.yaml} & Deployment + Service & Env depuis Secret/ConfigMap, probes /health, ServiceMonitor \\
    \texttt{web.yaml} & Deployment + Service & nginx, port 80 \\
    \texttt{secret.yaml} & Secret & JWT, mot de passe PG, DATABASE\_URL \\
    \texttt{configmap.yaml} & ConfigMap & CORS \\
    \texttt{ingress.yaml} & Ingress (×2) & web et API sur hôtes distincts, TLS \\
    \bottomrule
  \end{tabular}
\end{table}

\subsection{Ingress et routage}

Deux ressources Ingress séparent frontend et API :
\begin{itemize}
  \item \texttt{pfa.elyes.dev} → Service \texttt{web} (port 80) ;
  \item \texttt{api.pfa.elyes.dev} → Service \texttt{api} (port 8000 en cluster uniquement, non exposé sur l'hôte).
\end{itemize}

Le frontend est construit avec \texttt{VITE\_API\_URL=https://api.pfa.elyes.dev} : les appels REST ciblent directement le sous-domaine API.

Les annotations \texttt{cert-manager.io/cluster-issuer: letsencrypt-prod} provisionnent les certificats TLS.

\subsection{Sondes et ServiceMonitor}

\begin{itemize}
  \item \texttt{startupProbe} et \texttt{livenessProbe} sur \texttt{/health} ;
  \item \texttt{ServiceMonitor} pour Prometheus si \texttt{monitoring.enabled}.
\end{itemize}

\section{ingress-nginx sur le VPS}

\texttt{infrastructure/terraform/values/ingress-nginx-vps.yaml} configure le contrôleur avec \texttt{service.type: LoadBalancer} (k3s ServiceLB lie 80/443 sur le nœud), sans \texttt{hostPort} (contrairement à l'ancienne configuration kind).

\section{Cycle de déploiement}

\begin{enumerate}
  \item Bootstrap unique : \texttt{terraform apply} (plateforme) ;
  \item À chaque push sur \texttt{main} : build images → GHCR → \texttt{helm upgrade} (application).
\end{enumerate}

\section{Sécurité et limites}

\begin{itemize}
  \item un seul nœud — pas de tolérance aux pannes matérielles ;
  \item secrets production dans GitHub Actions et Kubernetes Secrets ;
  \item ports hôte 8000, 6001, 6002 non utilisés (réservés ailleurs sur d'autres serveurs).
\end{itemize}
